\section{Further remarks and questions}
\label{sec:remarks}

We conclude with several distinctions and open questions
arising from the results established above. The emphasis
is on identifying what the present argument does and does
not imply.

\subsection{The exterior class and ordinary topological complexity}

The ordinary tensor-valued Costa--Farber canonical class
$\mathfrak v^{2n}$ is an obstruction to motion planning
\cite{CostaFarber2010}. The exterior class
$\mathfrak a_{2n}$ is obtained from it by a quotient of
coefficient modules. Theorem~\ref{thm:intro-estimate}
shows that this exterior image is nonzero for closed
oriented negatively curved manifolds. This is useful
because the exterior coefficient module is a quotient
of the oriented simplicial chains of the full simplex.

However, the exterior construction does not give a
new value of $\TC(M)$. The equality $\TC(M)=2n+1$
was already known \cite{Dranishnikov2020}.
The additional content here is a quantitative
connection with multiplicity and a localization
formula whose terms are detected by
minimal-length magnitude homology.

\subsection{Comparison with cut-locus methods}

Geodesic complexity is sensitive to the structure
of the total cut locus. Existing work provides
upper bounds from decompositions of cut loci
\cite{RecioMitter2021,MescherStegemeyer2023}.
Our regularity condition instead controls all
simultaneously minimizing length branches,
including the highest multiplicity stratum.
The selection construction in
Section~\ref{sec:generic} uses this transverse
stratification to obtain an upper bound.

There is no general inequality identifying
geodesic complexity with the maximum number
of minimizing geodesics. Their equality in
Theorem~\ref{thm:generic-GC} is a consequence
of the specific geometric and cohomological
hypotheses, not a formal property of
geodesic complexity.

\subsection{The strength of the seminorm}

The seminorm $\|\cdot\|_\Delta$ is defined
by a concrete generating set in the
exterior augmentation coinvariants.
Our estimates show
\[
 c_n\|M\|^2
 \leq\|A_M\|_\Delta
 \leq\|M\times M\|.
\]
The inequalities need not be equalities.
Determining the exact value of
$\|A_M\|_\Delta$ even for closed hyperbolic
surfaces appears to require a more precise
description of bounded alternating
cochains on $\Gamma$.

The constants $c_n$ arise from the
norm-controlled lifting procedure.
They are not expected to be sharp in
general. An improvement would follow
from a smaller lifting norm for classes
built from the orientation cocycle,
without requiring a general isometric
isomorphism in bounded cohomology.

\subsection{Nongeneric metrics}

For a regular metric the highest
multiplicity stratum is finite and each
local contribution to the localization
formula has absolute coefficient one.
For a nongeneric metric one may have
isolated points with more than $2n+1$
minimizing geodesics, or positive-dimensional
sets of degenerate configurations.

The nerve construction remains available
for any negatively curved metric, but
the localization as a finite signed sum
of top-stratum points uses regularity.
A possible extension is to use relative
homology of the local simplex and its
lower-dimensional skeleton to define
local indices at isolated degenerate
configurations. Such an extension is
not proved here.

\subsection{Magnitude homology beyond degree two}

The identification
\[
 \MH_{2,d_M(x,y)}(M;x,y)\cong I(S_g(u,v))
\]
uses two special features: the length
is exactly the endpoint distance,
and the homological degree is two.
Higher magnitude homology groups
incorporate more complicated chains
and need not be determined only by
the cardinality of the minimizing
geodesic set. Their relationship
with higher-dimensional strata of
the minimizing-geodesic nerve is
a separate question.

The comparison with interval-poset
models in \cite{KanetaYoshinaga2021}
is a natural starting point, but
the existence of a global sheaf or
derived-category refinement is not
asserted in this paper.
